
\section{Literature Review and Research Gap}
\label{sec:lit}

Three bodies of work bear on the problem. The first supplies the routing
mechanisms and what is known about when they help. The second supplies the
objectives beyond travel time and the machinery for comparing them. The third
supplies the decision-theoretic apparatus for choosing among mechanisms and for
judging whether the choice generalises. Each is summarised for what it
establishes; Section~\ref{sec:gap} states what none of them settles.

\subsection{Dynamic and Traffic-Aware Routing}
\label{sec:lit:routing}

Route guidance begins from the observation that individually optimal and
system-optimal assignments differ \citep{wardrop1952}. Dynamic guidance replaces
static link attributes with observed or predicted conditions, and proactive
formulations spread vehicles over admissible alternatives rather than
concentrating them on the current fastest path; the flow-balanced $k$-shortest
path construction of \citet{pan2013proactive} is the direct ancestor of the
capacity-aware policy used here.

A second and more cautionary strand establishes that guidance is not
monotonically beneficial. Providing travellers with information need not reduce
congestion and can increase it, both in equilibrium models
\citep{arnott1991information} and in networks where added capacity or added
information degrades outcomes \citep{acemoglu2018braess}. The size and even the
sign of the effect depends on how many vehicles are guided: \citet{yang1998penetration}
shows that benefits vary non-monotonically with market penetration, and recent
work reaches the same conclusion from other directions. \citet{toso2024detrimental}
show in a dynamical network-flow model that recommendation-following can produce
partial demand transfer and congestion build-up whose severity is governed by the
penetration rate, and \citet{cornacchia2026concentration}, simulating commercial
navigation services on three Italian cities, find that route diversity falls as
adoption rises and that emission benefits diminish, vanish or reverse beyond a
city- and service-specific adoption threshold.

This literature therefore establishes that a routing policy's value is
state-dependent and penetration-dependent, and that the dependence can change
sign. What it does not do is treat the choice among established policies as the
object of study: it compares mechanisms, or studies one mechanism as adoption
varies, rather than asking which of several should be active in a given
operating condition and what that choice is worth.

\subsection{Environmental, Multi-Criteria and Context-Dependent Routing}
\label{sec:lit:multi}

Routing on criteria other than time is well established. Eco-routing assigns
paths using modelled emissions or fuel consumption
\citep{boriboonsomsin2012eco}, and microsimulation supplies per-vehicle emission
models that make such criteria computable inside a simulation
\citep{krajzewicz2015emissions}. Reliability-aware guidance treats the variance
of travel time as a cost alongside its mean; the mean--variance formulation of
\citet{sen2001meanvariance} is the ancestor of the reliability-aware policy used
here.

Where several criteria are in play, the question of which policy is
``preferred'' becomes a question about the criterion vector rather than about a
scalar. That distinction matters here for a specific reason: a policy can be
non-dominated on the criterion vector while never being best on any single
criterion, and a study that scalarises early will not see it. This study
therefore reports Pareto structure directly, and uses scalarisation only as a
declared sensitivity analysis. It does not propose a multi-criteria decision
method, and it deliberately avoids monetary scalarisation, because no value of
time, value of reliability or carbon price was available for this setting that
could be sourced rather than assumed.

\subsection{Algorithm Selection and Decision-Oriented Evaluation}
\label{sec:lit:as}

The selection problem was formalised by \citet{rice1976} as a mapping from
measurable properties of a problem instance to the algorithm expected to solve
it best. Solver portfolios put that mapping into practice
\citep{xu2008satzilla}, benchmark libraries fix how the resulting systems are
described and compared \citep{bischl2016aslib}, and surveys together with
instance-space analysis set out the methodology and where it breaks
\citep{kotthoff2014survey,smithmiles2009,smithmiles2014instance}. Two reference
points from that literature are used throughout this paper: the single best
solver, the strongest deployable fixed choice, and the virtual best solver, the
per-instance best chosen with hindsight, which is a bound and not a method.

The same literature supplies the argument for evaluating a selector by the cost
of its decisions rather than by its predictive accuracy. A model may predict
outcomes poorly and still decide well if it preserves the ordering of the
alternatives, and may predict well and decide badly if a small error falls near a
boundary \citep{elmachtoub2022spo}. Where a selector must be trusted outside the
conditions it was fitted on, the relevant machinery is that of distribution shift
and selective prediction: the failure of exchangeability under covariate shift
\citep{shimodaira2000}, conformal prediction and its weighted extension
\citep{vovk2005alrw,lei2018conformal,tibshirani2019covshift}, and the classical
reject option \citep{chow1970reject}. Gradient boosting supplies the regression
model used here \citep{friedman2001gbm,ke2017lightgbm}.

\subsection{Research Gap and Positioning}
\label{sec:gap}

Each strand answers part of the question and none answers the whole of it.
Traffic-aware routing research demonstrates state dependence but keeps the
mechanism, not the choice between mechanisms, as its object. Multi-criteria
routing research supplies objectives and the means of comparing them without
asking whether, in a particular network, those objectives separate the
candidates at all. Algorithm-selection research supplies baselines, regret and
generalisation testing, but its traffic applications are scarce and none pairs a
controlled factorial over operating conditions with a shift analysis reported by
shift type.

The gap is accordingly not a claim that regime-dependent routing is unstudied;
it plainly is not. The gap is that the three quantities named at the end of
Section~\ref{sec:intro} have not been measured on a common benchmark, so their
relationship is unknown, and a reader has no basis for converting evidence of the
first into an expectation about the third. Every ingredient used below has
precedent and is cited above. The contribution is the joint measurement, on a
design large enough to separate the three quantities from one another and from
the replication noise of the simulator, and the reporting of what that
measurement returns, including where it returns nothing.
